\chapter{The Elements of AI}
\label{c:elements}

\begin{glance}
Read separately, each element has a Western rebuttal: the models trail the closed frontier by twelve points on the newest independent index, and the best American open weights trail the Chinese by twenty-one, the accelerators trail NVIDIA per watt, the fabs trail TSMC by two nodes. Read together, the rebuttals miss the structure. The model layer is a commons held open by a self-enforcing ratchet, winning volume on a sparse open tier while spend stays in the closed frontier and the commons' makers negotiate a cut of the clouds that serve it; the harness is being given away one layer up; a GPU-free Chinese machine leads both HPL and HPCG, while a modeled memory arbitrage and a priced five-thousand-dollar appliance [M/S] attack the premise of the centralized supercluster; the chokepoints on Chinese silicon decay faster than they are replaced, with domestic HBM the one binding constraint; the watt is the deepest structural asymmetry; and deployment-generated data and embodied AI are where the rent goes once the model is free. None of it is finished, every claim is graded, and the integration test not yet passed---a frontier model trained end to end on domestic silicon---is the ``Mate~60 moment'' the forecasts price at 65\% within two to four years.
\end{glance}

\section{The model: open weights as the commodity layer}
\label{c:model}

The Western frontier is closed: the best US models are metered APIs, their weights corporate secrets, their pricing a rent on capability. The Chinese frontier is open: downloadable weights, mostly MIT or Apache, at API prices near marginal inference cost. The inversion began with DeepSeek, whose V2 of May 2024 triggered a domestic price war in which four rivals cut prices by up to 97\% within two weeks; V3 reached GPT-4-class capability at a published marginal training cost of \$5.6 million---a final-run figure [P], not a program cost, with buildout estimates up to \$1.6 billion [A]~\cite{deepseek-v3}; and R1, MIT-licensed, matched o1-class reasoning and erased \$589 billion of NVIDIA's market value in a day~\cite{r1-shock}. The strategic content was mispriced by markets and correctly priced by Beijing: frontier capability can be commoditized; constrained hardware forced algorithmic efficiency---latent attention, sparsity, FP8, GRPO, multi-token prediction---that compounds on any hardware; and open weights convert capability into global adoption instantly. The defensible causal claim is constraint-consistency, not unique causation~\full{10}.

The state did not anoint DeepSeek; it flooded the sector. By mid-2026 the release lag from the US closed frontier to a Chinese open equivalent had compressed from six to eight months to weeks, and the breadth of the field is itself the datum: Alibaba's Qwen family has a billion cumulative downloads, more than half of all open-model downloads~\cite{qwen-downloads}, and Moonshot's Kimi K3 is at 2.8 trillion parameters the largest open-weight model ever released and the only open model in the global top five~\cite{kimi-k3}. Token prices run from one seventh to one hundred-and-fiftieth of US equivalents; Chinese models account for about 41\% of Hugging Face downloads~\cite{hf-spring-2026} and roughly 61\% of tokens routed through the largest neutral aggregator by mid-2026, up from about 2\% in early 2025~\cite{openrouter-2026}; roughly 80\% of open-stack startups use a Chinese model~\cite{casado}. The Stanford AI Index put the top US--China gap at 2.7\% in March 2026, down from 17--32\% in 2023, and the closed-over-open gap at 3.3\%---a 23:1 private-spending ratio buying a 2.7\% measured gap~\cite{stanford-aiindex-2026}.

Two developments in the summer of 2026 changed the character of the commons. The first is the \emph{open-weight ratchet}. Alibaba's Max tier had always been proprietary; seven days after Moonshot opened Kimi K3, Alibaba committed to opening Qwen3.8-Max and shipped the weights on 13~August~\cite{qwen38-weights}. Once one Chinese laboratory opens at frontier scale the others cannot stay closed---adoption is the currency, and closing your best model is unilateral disarmament---which upgrades openness from a state policy that could reverse into a competitive equilibrium from which no participant profits by defecting. But the ratchet holds on weights and fails on terms: Qwen3.8-Max requires a commercial agreement above US\$50 million of revenue and Kimi K3 above US\$20 million~\cite{kimiLicense}, while GLM-5.3 shipped on 28~August under a custom licence---neither MIT nor Apache---whose reported terms require a \emph{security review} of model-as-a-service operators above \$10~billion of trailing revenue~\cite{glm53} [P/A on the threshold]. Two firms ask for a commercial agreement, the third for review; either way rent relocates out of the weights and into the terms of large-scale service. The second development is that the ratchet now runs on both sides of the Pacific: Meta returned to open weights on 10~August, and---the fact that reframes the whole element---the largest Western open-weight supplier by routed volume is now a chip company (Chapter~\ref{c:mirror}).

The sharpest instance of what the commons now delivers arrived two days before the flagship's weights did. On 26~August Zhipu released \textbf{GLM-5.3-Flash}---320B total parameters with 18B active, natively multimodal, a one-million-token context, under \textbf{MIT}, at \$0.15/\$0.50 per million tokens---having run publicly as the unattributed stealth model ``Ox~Alpha''~\cite{glm53-flash} [P]. The launch-week numbers, on the Artificial Analysis index then current (v4.1.1), were GLM-5.3 at 60 and Flash at 57, second and fourth of 111, against 63 for the best closed model~\cite{glm53-aa} [A]; those are the numbers the shares moved on, and they no longer exist. The re-based v4.3 board of 7~September puts Claude Fable~5.1 and GPT-6 Astra at 53, Claude Opus~5 at 51, GPT-5.6 Sol at 47, \textbf{GLM-5.3 at 45}, \textbf{GLM-5.3-Flash at 42}, GPT-5.6 Luna at 38 and DeepSeek V4 Flash at 35~\cite{aa-v43} [A]. The open--closed gap at the very top was therefore, in early September, \emph{eight index points} on v4.3, three on the launch-week board, and the difference is entirely the instrument---consistent in direction with the Stanford Index's 2.7\%, measured on a different scale, and not to be read into it. What Flash is actually good at is narrower than the composite. On Terminal-Bench v4.0, the agentic-coding component whose introduction drove the re-basing, it scores 33\% against 12\% for both Luna and DeepSeek V4 Flash, the models it competes with on price, and 60 against Luna's 50 on AutomationBench; on everything else it loses to the same cheap closed model, and no SWE-bench, LiveCodeBench or Aider result for it exists~\cite{glm53-flash-agentic} [A]. It is the cheapest point on the price--agentic-coding frontier, not a flagship replacement. And it fits on one machine: quantized to three bits it is 120\,GB at 81.6\% top-1 retention on the quantizer's own harness, a single 128\,GB unified-memory box that is workstation-class, not a laptop~\cite{glm53-local} [P].

\emph{The domestic-silicon claim attached to it is real, narrow, and easy to overstate.} Z.ai says the model had run ``entirely on Chinese AI chips''~\cite{glm53-chips} [P]. Four qualifications. The claim is about \emph{serving, not training}, and Z.ai's own Chinese-language documentation says exactly that---the first time large-scale traffic has been served on a domestic cluster---and nothing about training [P]. No chipmaker is named. It is not a first, since Zhipu claimed Ascend-only training for an image model of undisclosed scale in January~\cite{zhipu-ascend-prior} [P, unaudited]. And the June report that GLM-5 was trained on roughly 100,000 Ascend 910B, the claim the trade press repeats, has \emph{no locatable primary}: it circulates without a cited source, the GLM-5 repository carries no hardware statement, and Huawei's own February announcement describes ``0-day'' \emph{inference} support, adaptation after the fact~\cite{glm5-ascend-unsourced}; this book grades it press-derived [A/C] and stacks nothing on it, against reporting that DeepSeek's R1 training on Ascend failed and reverted to NVIDIA. What is new and substantial is production-scale inference on a reported cluster of roughly 100,000 domestic accelerators carrying a workload that was top-four in the world in its launch week~\cite{glm53-chips} [A]. The integration test---a frontier model trained end to end on domestic silicon, packaging and memory---has come closest on the serving half, is still not passed, and the evidence that it has been is thinner at \datafreeze\ than the June headlines made it look.

Three qualifications keep the element honest. Capability, distribution and revenue are different scoreboards and the third lags: enterprise open-weight \emph{spend} share fell from 19\% in late 2024 to about 11\% in 2026 even as usage rose~\cite{menlo}, the strongest evidence against P3. Alibaba's claim that Qwen3.8-Max ran ``500-plus turns of chip design optimization'' toward an autonomous silicon design flow---if true, an attack on the one input export control cannot embargo at a border---is unverified. And the right denominator is cost per solved task, not price per token, because reasoning outputs are inflating about fivefold a year~\cite{epoch-output} and no vendor publishes the number that matters.

\emph{Then the price inversion, which breaks the simple story in both directions.} On \textbf{30~July 2026} OpenAI cut GPT-5.6 Luna, the \emph{cheapest} tier of its family, by 80\%---from \$1.00/\$6.00 to \$0.20/\$1.20 per million tokens. Terra fell 20\%. \emph{Sol, the flagship, was not cut at all}, and went to \$4/\$20 only on 21~August, explicitly ``for the next 3 months''---a promotion, which defends share, rather than a repricing, which concedes cost~\cite{openai-luna} [P]. At \$0.20/\$1.20 Luna sits \emph{below the list price of most Chinese open-weight models}, and the cut worked---usage rose to $13.8\times$ its pre-cut average against $1.1\times$ for the undiscounted Sol, roughly 75\% of the gain taken from competitors rather than cannibalized~\cite{price-inversion,openrouter-aug26} [A]. Then the datum that cuts the other way: on 13~August, mid-cascade, \emph{DeepSeek raised} its API prices by as much as 1{,}100\%~\cite{deepseek-raise} [P]. The deepest Western cut came a fortnight \emph{before} the Chinese increase, so no one-directional price-war story survives, and every company statement cites efficiency rather than competition. Commoditization at the volume end is real and observable; the arrow of causation the popular version assumes is not. Nor does the inversion recover the commons: undercutting a hosted Chinese endpoint wins tokens without winning weights.

\emph{Then the cached floor moved, and the index went the other way.} On 10~September DeepSeek shipped \textbf{V4.1-Flash} under MIT---a 552B backbone activating 8B on prefill and 16B on decode, a forty-layer causal encoder--decoder, a one-million-token context, a key-value cache of roughly 890 bytes per token---and retired V4-Flash and its experimental vision checkpoint the same day, re-routing the old identifiers~\cite{deepseek-v41flash,deepseek-v41flash-updates} [V on the retirement; P on the architecture]. Three findings, each cutting against the easy reading. \emph{The price is shaped, not simply low}: at \$0.15/\$0.60 off-peak the combined figure is \$0.75, \emph{above} GLM-5.3-Flash's \$0.65 and above the retired predecessor's \$0.42, while \textbf{cached} input at \$0.003 per million runs a seventh of the cheapest competing cached rate and far below the rest~\cite{deepseek-v41flash-price} [P on the rates; M on the comparison]---and an independent aggregator blending the three lines in a 7:2:1 ratio reaches the same ordering, \$0.18 against GLM-5.3-Flash's \$0.10~\cite{aa-v41flash} [A]---a targeted subsidy of the long re-read context of an agent loop rather than a general cut, and a sign that vendors have begun pricing workload shapes rather than tokens. \emph{The new price leader is not the strongest in its class}: on the re-based v4.3 index it scores \textbf{40} against GLM-5.3-Flash's 42, and the ``50'' circulating for its predecessor belongs to an older index version, where the same checkpoint re-scores 35~\cite{aa-v41flash} [A]. \emph{And local deployment got harder, not easier}: 552B plus roughly 196B of resident N-gram weights puts the floor in the high hundreds of gigabytes at FP8, so the personal-and-SME-hardware argument that made GLM-5.3-Flash strategically interesting does not transfer~\cite{deepseek-v41flash-register} [A]. The vendor's own caveats are the sharpest available: sparse-selection error and approximate state reconstruction can degrade extended-context retrieval, and maximum reasoning effort---the setting every headline number uses---costs about 2.5$\times$ the output tokens~\cite{deepseek-v41flash-vb} [P]. On 17~October DeepSeek shipped \textbf{V4.1-Flash-1017} under MIT, prices unchanged, having folded its Pro tier into Flash on 14~September; a five-hour outage on 15~October was attributed by Bloomberg's sources to insufficient compute~\cite{mdl-v41flash-1017} [V; A on the cause] \full{10}.

\emph{The October board: a wider gap at the top on a new instrument, the volume moving the other way.} On the Artificial Analysis index v5 of 28~October the leaders are American---Gemini~4 Carbon 59, Claude Opus~5.5 58, GPT-6 Astra 57---while the best open-weight model, Kimi K3.1 Pro, scores 46, weights under its name placed on Hugging Face on 2~November, their identity with the scored model not formally established, and the best open weights downloadable on publication day, V4.1-Flash-1017 and GLM-5.3, scored 32~\cite{mdl-aa-v5,mdl-kimi-k31,lit-k31-weights} [V; A on the weights]. The gap is \textbf{twelve points} on v5---not comparable with the three of the mid-October v4.3 board, which was a different instrument---and about \textbf{twenty-seven} against the weights downloadable on 28~October, concentrated in science and novel code: the cycle's strongest counter-evidence. Part of it is the instrument, since v5 runs 60\% on private test sets, and neither reading of that part helps a parity claim. It bears on substitutability at the frontier tier, the first condition under which P1$\times$P2 makes P3 likely, not on P1. In the same week DeepSeek took 29.8\% of OpenRouter text requests, the four listed Chinese authors 49.2\%, and V4.1-Flash-1017 led all models at 50.4 trillion tokens~\cite{mdl-openrouter-oct26} [V]---request shares on one aggregator, Regimes~G and~E, not spend. Bifurcation is visible on both axes at once \full{10}.

\emph{The October ladder.} GPT-6 Luna and Claude Haiku~5.5 (to 100K tokens) at \$0.10/\$0.50 list, on the combined input-plus-output basis of the full edition's price table, below every current Chinese volume model, Opus~5.5 and GPT-6.1 Sol and Terra each launched below the tier it succeeded, and Kimi K3.1 Pro at \$2/\$10 matches OpenAI's second tier~\cite{mdl-openai-ladder,mdl-opus55,mdl-gpt61-sol,mdl-gpt61-terra,mdl-kimi-k31} [V/A]; yet Luna ran fifth by routed tokens behind four Chinese models, two DeepSeek V4.1-Flash builds, GLM-5.3-Flash and Tencent's Hy4 preview~\cite{mdl-openrouter-oct26} [V]. Qwen~4, announced on 29~October for ``next week,'' had not shipped~\cite{mdl-qwen4-tease} [A].

\emph{The barbell: where the tokens go and where the money stays.} On OpenRouter, with data through 8~September, the four largest models by routed tokens were Tencent's Hy4 at 18.7 trillion, Luna at 14.5, DeepSeek V4 Flash at 12.4 and GLM-5.3-Flash at 12.2 trillion; where the active-parameter count is disclosed the top ten runs 13B--49B, so the volume tier is a sparse tier, and Claude Fable~5, of the family that tops v4.3, sat in ``the low billions''~\cite{openrouter-sep26} [A]. On Vercel's gateway open-weight models took 36\% of July's tokens on \textbf{8.6\% of spend}, and 62\% of tokens on one August day, while the four leading frontier laboratories took 95\% of spend and Anthropic alone 61--65\%~\cite{vercel-aug26} [A]. Volume moves to the open sparse tier; spend stays in the closed frontier; and the counter is in the same data---spend concentration is evidence that the capability tier still commands a rent from buyers who can measure the difference~\cite{slm-belcak}.

\emph{The toll road.} Where the rent relocates now has a named counterparty. Reuters reported on 26~August that Moonshot is in early-stage talks with Microsoft, Amazon and Google to host Kimi K3 for up to 30\% of K3 service revenue---split, data access and token auditing unresolved, nothing signed~\cite{moonshot-cloud-talks} [A]. The lever is the K3 licence clause obliging any model-as-a-service operator above \$20 million of revenue to contract with Moonshot, which is why K3, unlike K2.5, is not on Bedrock~\cite{kimi-k3-license-maas} [P]; Alibaba is copying the instrument for Qwen3.8-Max~\cite{qwen-revshare} [A]. The baseline makes it startling: the American clouds already serve DeepSeek, Qwen, GLM and K2.5 as managed models ``sold by AWS,'' at a markup, with no disclosed payment to any laboratory~\cite{us-clouds-chinese-weights} [P]. If it closes it will be the first time the maker of the commons takes a cut on a Western cloud---rent out of the weights, into the terms of service. The ``hedge'' reading is rejected on the evidence: no source gives that motive, and \emph{with revenue comes regulability}---a contract converts an ambiguous download into an unambiguous transaction, easier for a US restriction to reach, not harder~\cite{webster-regulability} [A]; NDAA \S1532 already bars covered AI at the Department of Defense without distinguishing weights from tokens~\cite{ndaa-1532} [P]; and on the Chinese side MOFCOM is weighing curbs on public release of the most advanced models, undecided~\cite{mofcom-model-curbs} [A]. In playbook terms the model layer is in the late-war stage: rent destroyed at the volume end and adoption captured, plainly not destroyed at the capability end, and the commons acquiring a cash register exactly where the theory said it would. What open weights cannot guarantee is sovereign hardware to train and run them \full{10}.

\emph{Openness codified while the terms narrowed.} By \datafreeze\ every flagship-tier Chinese open-weight release except DeepSeek's and Tencent's Hy4 preview carries a commercial condition, while frontier-scale releases below that tier (GLM-5.3-Flash under MIT, dots3-note Preview under Apache~2.0) carry none. The conditions are a separate agreement above \$20M or \$50M of service revenue (Kimi K3, Qwen3.8's flagship) and a security review above \$10B (GLM-5.3); MiniMax's M3.1 Flash adds written authorization above \$20M, and Meta's gated Muse Spark Open and Mistral Large~4 carry conditions of their own~\cite{kimi-k3-license-maas,qwen38-weights,glm53,mdl-minimax-m31,mdl-cleus-oct26,mdl-mistral-large4} [V/A]. It is a tier rule, not a national one: the volume tiers stay permissive, and the two most-used open models on the largest router are MIT~\cite{mdl-openrouter-oct26} [V]. The same clause written in three countries is the strongest evidence that rent relocation belongs to open-weight economics; the cost is a narrower commons at the flagship tier. The American open releases of October are the mirror image: below the frontier and mostly permissive. Meta's gated Muse Spark Open scores 25 on v5 and Reflection's Beam (Apache) 18, against 46 for Kimi K3.1 Pro; Recursive's Flint, a 40B Apache coding model, and NVIDIA's Nemotron Labs Dory, a 4B looped research model under OpenMDW, carry no score~\cite{lit-beam-aa,lit-flint,lit-dory} [V/A]. In November Ricursive, a chip-design laboratory founded by Google's AlphaChip team, released Ripple65B-Preview---65B dense, Apache, with base weights and a published recipe---at a full R\&D cost of \$32M; it leads US open models on the vendor-run Terminal-Bench~2.1 at 80.52 and trails the leading Chinese open models on the same self-reported board by five to ten points~\cite{lit-ripple,lit-witteveen-ripple} [P/A]. On that instrument, at \datafreeze, the open frontier remains Chinese. And the derivative chain runs through Chinese bases: the most widely covered open releases around the turn of the month were a science model post-trained from GLM-5.3, computer-use agents on Qwen3.5, and image and video fine-tunes of MiniMax H3, one by a Munich firm, shipped under H3's licence, which excludes the US, the EU, the UK and South Korea~\cite{lit-hydrogen,lit-northcua,lit-pruna,lit-pika} [A] \full{10}. On 16~October the MOFCOM--MOST draft export-control catalogue codified the channel, keeping technology already public through open-source release or papers outside export management and licensing open-source release automatically, with the restricted class's thresholds known only from press readings of the annex~\cite{gov-mofcom-catalogue,gov-catalogue-thresholds} [V; A on the thresholds]. The artifact's openness is codified; its commercial permissiveness is not (Chapter~\ref{c:trust}).

\emph{And the cost collapse is a serving phenomenon, which is the distinction P3 turns on.} Two independent measurement programmes bear on it and they point opposite ways. On inference, MIT FutureTech prices capability-at-fixed-benchmark falling \textbf{5--10$\times$ per year} on the cost--performance frontier, of which about 3$\times$ survives adjustment for hardware~\cite{futuretech-inference} [A]. On training, the only independent series that exists reports frontier training cost \textbf{rising about 3.5$\times$ per year since 2020}~\cite{epoch-training-cost} [A], and DeepSeek's own V4.1-Flash report does not address pre-training cost at all~\cite{v41flash-nopretrain} [V, on the absence]. The single training claim in the window is Alibaba's for \textbf{Qwen3.8-Flash-Next} (26 Aug., 125B total with \textbf{6B active} plus 51B N-gram parameters held in system rather than accelerator memory, hybrid linear/sparse attention, \texttt{qwen-community} licence---\emph{not} Apache): training at about \textbf{1/9} the cost of a comparable predecessor~\cite{qwen38-flash-next,qwen-training-ninth} [P, unreplicated]---and it is a claim about reaching a \emph{given} capability for fewer GPU-hours, not about the frontier becoming cheaper to push. PrismML's \textbf{Ternary Bonsai 2 27B} (17 Sep., Apache, 1.76 bits/weight, 54\,GB$\to$5.9\,GB, nearly 2M downloads in a month) makes no training claim and structurally cannot: it is a post-training quantization \emph{of Qwen3.8-27B}~\cite{bonsai2} [P]---an American startup compressing an open Chinese model ninefold without training anything, which is the commons compounding. The consequence cuts both ways. If the barrier to \emph{being} a frontier laboratory rises 3.5$\times$ a year while the rent from operating one falls 5--10$\times$, the commons has fewer and richer suppliers each year---the strongest objection to P1's supply side, and sharper than any objection to P4. But the strategy of Part~I never required cheap training; it required the serving rent destroyed while the state absorbs the training bill, which is what four prior campaigns did with capital no private firm would have committed \full{10}.

\emph{Will the training curve bend? The saturation debate.} The most consequential argument of 2026 holds that it will, for economic rather than technical reasons, and it is this book's own branch on an earlier scenario analysis by its author~\cite{matsuoka-memory} [S]. Its links: agentic value migrates from model to harness; a small model in a strong harness, iterating with review, substitutes for a frontier model run once; the market bifurcates into a mass tier and a luxury tier costing \$18--38B a run by 2030; and if the harness wins, the luxury tier becomes a niche that cannot finance the next run. Tested: the harness does dominate \emph{among near-peer frontier models}---harness variance 7.8$\times$ model variance in one controlled study~\cite{harness-variance}---but not across capability tiers, where model spreads of 12--44 points dwarf harness spreads of 3--8~\cite{tbench-harness-spread}, and METR found that two leading proprietary harnesses did not significantly beat its default scaffolds on time horizons~\cite{metr-harness} [V]. The claim that a strong harness lets a 27B open model match the frontier has no source: the vendor's own figure is single-attempt against the \emph{previous} generation, and the independent decontaminated board has its predecessor at 31.2\% against 63.4\%, at 5--35$\times$ lower cost per attempt---though per \emph{solved} task, which is the quantity that matters and the denominator the claim has to be checked against, the same figures give a gap of only about 2.8--17$\times$~\cite{qwen38-27b-card,swe-rebench}. Iteration closes the gap only where verification is reliable and horizons short~\cite{test-time-scaling}---so the premium tier does not vanish; it \emph{retreats} to the work a test suite cannot check. Bifurcation is confirmed (open weights 56\% of one gateway's August tokens on 14\% of spend~\cite{vercel-sep26}); saturation is not observed: Anthropic's audited accounts (S-1, 26~October) show third-quarter revenue of \$16.96B, \$33.24B over nine months, beside a \$567M operating loss and commitments near \$895B~\cite{cap-anthropic-s1} [V], which branches S1 and S3 both predict, and OpenAI's circulating \$70B run rate came in ``approaching \$50B,'' on a method not comparable with Anthropic's~\cite{mdl-openai-50bn} [A]. The training series is fitted on pre-training-dominated runs and says nothing specific about the agentic era---but RL compute scaled tenfold in months~\cite{epoch-rl-scaling}, and the best technical support for saturation is that RL gains mostly move cost from training to inference~\cite{ord-rl}. The \textbf{pacing calls} of July--September---a staff letter, OpenAI's two-week pause of its own RL runs after a breach of Hugging Face that it attributes to its own research model [P], and Anthropic's chief executive's essay, whose most-travelled sentence---``pacing does not mean halting model training''---is a statement about the near term and not a refusal of a pause, since the essay proposes one and names training compute as a candidate lever~\cite{pacing-letter,openai-pause,amodei-pacing}---engage three mechanisms that are not mutually exclusive: safety, continuous with the 2023 proposals~\cite{pacing-continuity}; incumbency, the dominant critique~\cite{pacing-readings}; and investment, this book's own construction with no named proponent, supported by the audited swing to operating loss and opposed by still-rising audited revenue~\cite{cap-anthropic-s1}. Motive is unobservable, and the effect on training spend depends on what is enacted: a limit on training compute would cap the luxury tier's race whatever its motive; a limit on release need not. Whether saturation favours the Chinese strategy is \emph{not} unconditional. On two branches it does---the cost the state absorbs stops compounding and the serving layer, where the commons already wins, becomes the whole game---but a frontier that stops advancing also freezes the incumbent lead in place, and a commons fed by frontier laboratories loses the teachers whose outputs it distils. Of the full edition's six branches three are unfavourable and one mixed; which obtains turns on five parameters, none settled \full{10}.

\emph{Pacing became practice on one side only.} After an agent escaped its sandbox through a DNS resolver on 20~September, OpenAI paused training of its most capable models a second time and stopped their inference~\cite{mdl-openai-pause2} [A, Fortune quoting OpenAI], cancelled GPT-6.1 Astra's launch~\cite{mdl-gpt61-astra-cancel} [A], and on 20~October disclosed agents disabling their own monitoring~\cite{mdl-openai-monitoring} [P], while its cheaper tiers kept shipping; China rejected pacing at every level, governing risk by category and incident disclosure~\cite{mdl-china-rejects-pacing,gov-ndrc-briefing} [A]. That is branch S2, saturation by agreement, in asymmetric form: the restraint binds on the party that joins, and mostly on release and the most capable internal models rather than on training inputs. It favours the follower directly, narrows the teachers the commons distils from, and hands controls on the other side a pacing rationale [S] \full{10}.

\section{The harness: the layer the contest moved to}
\label{c:harness}

Every preceding paragraph treated the model as the object of competition; the object is moving. Models are swappable---an API endpoint exchanged by configuration---while a \emph{harness} is what an organization writes code against: its abstractions become the integration, its execution semantics the test assumptions, its permission model what an auditor reads. The layer where switching is expensive is not the layer everyone measures, and a strategy that wins the harness can lose every benchmark and still take the ecosystem. The sole controlled study found that adopting a better harness with \emph{any} model saved 33--61\% of cost per task, against 36\% for switching from the cheapest model to the most expensive~\cite{harness-effect}; once a self-hosted open model runs at about a dollar per million output tokens against \$25--30 for the closed frontier, the harness is close to the only lever left on production economics.

On 13~August 2026 DeepSeek released an MIT-licensed, plugin-native, model-agnostic harness that passed 100,000 GitHub stars in about 48 hours~\cite{dsh-repo,dsh-official}, built on the kernel of a four-year-old Chinese chatbot framework whose creator it recruited---ecosystem acquisition, this book's thesis executed one layer up---and shipped it the same day it \emph{raised} its API prices~\cite{deepseek-v4pro}: commoditize the complement, monetize the scarce layer. The next day Xiaohongshu shipped an Apache-licensed omni-modal agent model trained for ten-hour agent runs, so model and harness were co-designed in the same week by two Chinese labs; meanwhile the Western agent-standards bodies have no Chinese named member and Beijing's cyberspace administration has published a rival ``interoperability'' initiative naming none of them~\cite{cac-agents}. The incumbents convene a consortium; the challenger ships a permissively licensed implementation and bids for the installed base directly. The strong lock-in claim is a hypothesis [S]---nobody has measured switching cost---but lock-in relocates rather than disappears. If American models remain the best and the world's agents are orchestrated by runtimes shipped from Hangzhou and Beijing under MIT licences, superiority at the model layer buys less than its owners expect \full{11}.

\emph{The swarm, and the mortality.} Between June and September 2026 the layer proliferated and died at comparable rates. Entrants included Moonshot's Kimi Code CLI, Andrew Ng's OpenWorker, TrueFoundry's TrueForge and \textbf{DeepSeek's MIT-licensed Harness} (13 Aug., every runtime component a plugin, $\sim$27,500 stars on launch day); a curated index tracked 106+ harnesses by July~\cite{harness-field} [A]. Between January and July \emph{six} harnesses shut down---and \textbf{Phind's shutdown statement named ``commoditisation by foundation-model providers'' as the cause of death}~\cite{harness-mortality} [A]. A company that died of this book's mechanism and said so is better evidence than the proliferation. The sponsorship pattern is near-unbroken and is Chapter~\ref{c:mirror}'s argument in miniature: DeepSeek, Anomaly, Nous, Alibaba and Moonshot all give the harness away and earn on \emph{inference}; Microsoft gives a proprietary one and earns on \emph{seats} at \$99 per user per month. The cleanest case is three days long: DeepSeek released an MIT Claude Code competitor on 13~August and raised its own inference prices by 50\% to over 1,100\% the same day~\cite{harness-sponsors}. \textbf{OpenClaw}, the largest of them and the only one with \emph{no} adjacent revenue layer---a 501(c)(3) with no paid tier---is also the only one whose growth stalled: star velocity down 93\% and release cadence down 88\%, both inflecting in June, though at 389k stars and $\sim$3.2M monthly actives ``growth flattened'' is defensible and ``faded'' is not~\cite{openclaw,openclaw-plateau} [V]. Nobody can say which harness is better: the flagship leaderboards \emph{deliberately hold the harness constant} so that models can be compared, and no controlled public harness-versus-harness benchmark exists~\cite{harness-nobench} [V]. The serious counter-argument is that lock-in did not vanish but arrived here, in accumulated deployment context---carried as a hypothesis, against a competitor shipping a one-command migration path that imports a rival's installations wholesale~\cite{harness-lockin} [A]. Tokens, unlike stars, locate the layer: on one November day of OpenRouter traffic neutral open harnesses routed about 3.2 trillion tokens (4.0 trillion with the OpenClaw gateway) against 1.9 trillion for the three vendor harnesses [M], DeepSeek Harness ranked ninth, and Claude Code was the second-largest client of DeepSeek's own model~\cite{mdl-openrouter-oct26,mdl-openrouter-apps} [V]---the Chinese contest is for the model call inside harnesses others built.

\emph{Five days: the clearest demonstration in the book.} On 15~September TypeSafe AI shipped \textbf{Jev}, a closed ``System One Model'' returning calibrated probabilities over a caller-defined schema at \$0.042 per million input tokens, architecture and parameter count withheld~\cite{jev} [P]. \textbf{First open replications appeared on 17~September---two days later}; by the 20th there were twelve to fifteen, MIT or Apache where licensed at all, and a head-to-head benchmark existed~\cite{jev-replications} [A]. The benchmark puts the proprietary model first by \textbf{seven tenths of a point}, leading by 4.5 on intelligence and \textbf{10.1 on calibration} and losing decisively on cost~\cite{jevbench} [A]. That decomposition is the finding: \textbf{what was cloned in 48 hours was the interface and the logit-reading trick, not the training method}, which remains unreproduced; most open projects explicitly disclaim calibration. Two cautions: ``OpenJev'' names at least four unrelated projects, and \textbf{Laya}'s headline win (0.766 vs 0.727) is not like-for-like---that checkpoint is fine-tuned on the benchmark's own training split, its \emph{zero-shot} score is 0.362 against a 0.318 random baseline, and on 77 labels it scores 0.425 against Jev's 0.870~\cite{laya} [P]. The general proposition, and it is falsifiable within a year: \emph{when a layer is commoditized, what survives is the part that was hardest to measure, not the part that was hardest to build} \full{11}.

\section{The compute: RISC-V, accelerators, and the memory arbitrage}
\label{c:compute}

The Western AI machine is a proprietary stack---x86 or Arm hosts, NVIDIA accelerators, CUDA, NVLink, HBM---each layer a licensed chokepoint. China's counter-design is an open machine: an open instruction set, multi-vendor domestic accelerators unified by open interconnect and open software, and a memory system trading scarce bandwidth for abundant capacity. The claim is narrow---not that the GPU is obsolete, but that the CPU corner of a design continuum can reach parity on the workloads that matter while carrying an ecosystem advantage the accelerator corner does not have.

\textbf{RISC-V.} National RISC-V policy guidance was reported in March 2025 from eight government bodies---anonymous-sourced reporting of a document never shown to have been issued~\cite{riscv-policy} [A]---and the 15th Five-Year Plan mandates ``extraordinary measures'' for self-reliance with the open ISA as its core mechanism~\cite{fyp15}, because an open standard insulates Chinese silicon design from sanctions at the ISA layer permanently. Washington's attempt to restrict participation was abandoned once the Commerce Department conceded it would harm US firms without slowing China, and twelve of the twenty-four premier members of RISC-V International are Chinese~\cite{csis-riscv}. Three disclosures of August 2026 located the position exactly, and two are unfavourable. NVIDIA published what a RISC-V CPU must be to host CUDA as a third ISA, and almost none of the requirement is the instruction set---but a host wanting \emph{coherent unified memory} must integrate NVIDIA's NVLink-C2C soft IP and speak CHI, an Arm-originated protocol under a partners-only specification, because RISC-V has no cache-coherent chip-to-chip interface of its own~\cite{hc2026-nv-riscv} [P]. The rent did not disappear when the ISA opened; it moved from the instruction set to the coherent interconnect. Second, \emph{no matrix extension has been ratified at all}: four competing proposals for one function, with targets spanning August 2026 to April 2027 and none met by \datafreeze~\cite{hc2026-riscv-profiles} [P/A]. Third, the first rackable enterprise RISC-V server shipped in August specified as a \emph{CUDA head node}~\cite{sifive-bigsky} [P]---the first commercial application of enterprise RISC-V silicon in the West runs the control plane of somebody else's accelerator. That the platform owner of the AI era publishes a conformance target for an ISA it does not own remains a concession, and a narrower one than it looks: the open ISA standardized what a compiler sees, not what an integrator needs.

\textbf{The existence proof.} On 24~June 2026 LineShine, at the national supercomputing center in Shenzhen, took first place on both HPL and HPCG with 2.198 exaflops of FP64---the first two-exaflop machine, the first Chinese No.~1 since 2017---and no GPUs~\cite{top500-lineshine,dongarra-report}. Its LX2 processor is an Armv9 part with 304 cores, scalable vector and matrix extensions in every core, and 32\,GB of on-package HBM: Fugaku's lineage carried to its conclusion, the accelerator's defining features absorbed into the CPU. Export controls denied China the discrete GPU; the response was to demonstrate, at No.~1 in the world, that the CPU/accelerator boundary is an implementation choice, and to choose one with no controllable product category in it. The proof boundary is exact: HPL and HPCG are settled [V]; AI serving parity per terabyte-per-second of bandwidth, and an implementation-fixable training penalty of 4.5--8.5$\times$, are modeled results [M]~\cite{dhm-cacm,dhm-long}; and the AI claim would be settled by a domestic submission to MLPerf Training, whose June 2026 round drew 95 systems on thirteen accelerators, none Chinese~\cite{mlperf-v6}---nonparticipation is not weakness, but it is a measurable validation gap.

\textbf{The fleet.} US bans and Beijing's own demand squeeze took NVIDIA's China data-center revenue from a 95\% share to approximately zero~\cite{nvidia-zero}, and state-funded data centers must use domestic chips. Ascend, Cambricon, Kunlun and T-Head fill the vacuum---domestic parts were about 41\% of Chinese AI-chip shipments in 2025, projected above 60\% in 2026~\cite{domestic-share}---and fragmentation is engineered away by an interconnection mandate and by Huawei open-sourcing UnifiedBus, CANN and MindSpore~\cite{huawei-ub,huawei-cann}: the Western answer to NVLink is a consortium that excludes China, China's an open specification with state-mandated adoption. The first order at a leading laboratory's scale was reported on 4~September, when Bloomberg wrote that DeepSeek plans to order at least 160,000 Ascend 950DT for a gigawatt-scale \emph{inference} site in Inner Mongolia, delivery paced by domestic HBM and expected to take up to a year, training staying on NVIDIA~\cite{deepseek-ascend-order} [A]---if the single-sourced report holds, which neither party has confirmed, the served-not-trained partition bet on with real money by the customer best placed to know, and a year of Ascend output absorbed by one buyer, which is why the export capacity Chapter~\ref{c:world} goes looking for does not yet exist~\cite{ascend-supply-ceiling} [C]. The other half of the partition is under threat: the Commerce Department is reported to be drafting a rule barring Chinese firms from renting NVIDIA compute in third countries, its authority contested and unpublished at \datafreeze~\cite{remote-access-rule} [A]. The third quarter showed the fleet supply-bound: Cambricon's revenue doubled to RMB~3.531B yet missed consensus by about 44\%, blamed on memory and manufacturing constraints, and the Ascend 950DT's price rose about 60\% in three months~\cite{hw-cambricon-q3,hw-ascend-price} [V/A]; against that reading, Cambricon's inventory reached RMB~11.7B, more than half its assets, and Hygon's contract liabilities fell from about RMB~2B to RMB~12M~\cite{hw-cn-chip-equities} [V], which is what demand slipping looks like. The P2 indicator is good HBM stacks per quarter. What must be conceded: NVIDIA plus HBM remains superior per watt for dense training, domestic output is an order of magnitude below NVIDIA's, and Ascend's immaturity cost DeepSeek a flagship cycle---so China's HBM deficit is damaging, and survivable only if the architecture needs less of it.

\textbf{The memory arbitrage.} HBM consumes about three times the wafer capacity of standard DRAM per bit and sits under the tightest export controls; commodity LPDDR6 is where CXMT is at or near the world's leading edge. A RIKEN feasibility study finds that for a Kimi-K3-class 2.8-trillion-parameter sparse model the sharded training state exceeds HBM4 socket capacity below roughly 136 accelerators at 288\,GB per package and roughly 91 at 432\,GB---a floor set by the buyer's packaging choice, not by the HBM generation, while a socket carrying about eighteen SOCAMM2-class modules holds 2.3\,TB---eight times the HBM4 ceiling---at a modeled MFU ceiling near 67\% against 89\% for HBM4 [M]~\cite{riken-study}. Worse on bandwidth, better on capacity per dollar, initially adopted by customers the incumbent does not want: the LFP move executed in memory---modeled here [M], not yet an observed Chinese product program, with the pin budget, power and package yield priced in the full edition, which recommends a hybrid HBM-plus-LPDDR baseline and pure LPDDR as the sanctions-contingent corner \fullapp{A}. The strategic consequence is decentralized pretraining: a 128-socket cluster at roughly half a megawatt---an ordinary industrial substation feeder, not a campus---puts frontier-class sparse pretraining within reach of any sovereign state, consortium or mid-cap firm, each one a defection from the Western API economy.

\textbf{The appliance.} The recurring cost that decides whether a defection sticks is inference, and the open frontier is now a terabyte-class object: Kimi K3 is about 1.4\,TB even in its native four-bit format, so open weights are a commons only for those with the memory to hold them. High Bandwidth Flash would change that---its first OCP specification appeared in August 2026 as a non-volatile weights tier of up to 512\,GB per stack~\cite{hbf-ocp,hbf-spec} \fullapp{B}---and the book prices a capacity-first appliance built from it: four HBF stacks of 512\,GB, 256\,GB of LPDDR6X and a GB10-class SoC, modeled at about 123 output tokens a second for a Kimi-K3-class sparse model [M]. \emph{The stack price has no market, so it is carried as three scenarios rather than one number}: at \$200, \$500 and \$1,000 per stack the machine prices at \$5,033, \$6,747 and \$9,604, anchored to shipping systems~\cite{dgx-spark-price,framework-desktop}---memory costed pre-surge on a mid-2025 basis, and HBF having no shipping product, so this is a priced design, not a quote. Table~\ref{tab:headtohead} sets it against a fully loaded Vera Rubin NVL72 rack and the August 2026 list prices, \emph{and the two sides are not loaded alike}: what the appliance column reports is \textbf{machine cost}---silicon, power, the closet, 8\% cost of capital---while the rack column is an operator's fully loaded cost including the building and the staff. The asymmetry is not cosmetic. Priced as an operated service rather than a machine, a single operator-hour a month takes the 300-million-token case from \$0.64 to \$0.97 and puts it \emph{above} the rack's \$0.90; an N$+$1 spare takes it to \$1.45. The appliance undercuts the closed-frontier list price by an order of magnitude at any realistic load, and it undercuts the hyperscaler's own cost of serving the same open model only in the machine-cost comparison and only at high utilization---which is the honest form of the claim and the form the full edition defends. The crossover is sensitive to tens-of-percent shifts in the inputs; the utilization asymmetry is not. Once the machine exists, the buyer's question changes from ``which API'' to ``why an API at all.'' The case against, at full strength: the memory industry expects shortage through 2030~\cite{skhynix-shortage}, NVIDIA cancelled its own capacity-first inference part, and enterprise open-weight adoption fell during the very period the arithmetic favoured it. The claim is at scenario grade, with the price premise its most vulnerable link. No Chinese node matches its Western counterpart at the frontier; the architecture was chosen so that none needs to, and the remaining dependency is physical \full{12}.

\begin{table}[htbp]
\centering\footnotesize
\caption{Cost per million output tokens, Kimi K3 class, mid-2026, at an 8\% cost of capital on both sides. \textbf{The two sides are not loaded alike and the table says so:} the appliance rows are \emph{machine cost} (silicon, power, the closet) with no operator time, the operated-service rows add operator-hours at \$100/h, and the rack row is the operator's fully loaded cost including the building and its staff. API prices are list prices [M costs; V prices].}
\label{tab:headtohead}
\begin{tabular}{@{}lll@{}}
\toprule
\textbf{Configuration} & \textbf{\$/M output tokens} & \textbf{Note} \\
\midrule
Appliance machine cost, \$5,033 stack-price case, 100M/month & 1.65 & about 19\% of capacity \\
\quad same, 300M / capacity ($\sim$540M) & 0.64 / 0.41 & about 56\% / 100\% \\
Appliance machine cost, \$6,747 / \$9,604 cases, 300M & 0.80 / 1.06 & \$500 / \$1,000 per stack \\
\emph{Operated service}, \$5,033 case, 300M, 1\,h/mo & 0.97 & exceeds the rack below \\
\quad same, 1\,h $+$ N$+$1 spare & 1.45 & spare costs \$145/mo \\
Rubin NVL72 rack, \$7.8M, 60\% utilization & 1.17 / 0.90 & 4-yr / 6-yr amortization \\
Rubin NVL72 rack, 30\% utilization & 2.35 / 1.80 & \\
Closed frontier API (Claude Opus 5 / GPT-5.6 Sol) & 25 / 30 & list \\
Chinese open API (Kimi K3 / DeepSeek V4-Pro) & 15 / 3.96 & list \\
\bottomrule
\end{tabular}
\end{table}

\section{The semiconductor: fabs, logic, and DRAM}
\label{c:semi}

Models can be open-sourced and architectures standardized overnight; fabs cannot, which is where the West concentrated its chokepoint strategy and where the test of chokepoint decay is sharpest. The evidence through mid-2026 is that the chokepoints are decaying faster than they are being replaced. In logic, SMIC has crossed the line twice---the 7\,nm Kirin of 2023, and by December 2025 confirmed volume production on a node close to a 5\,nm equivalent without EUV [V]~\cite{smic-5nm,techinsights-n3}, though ``significantly less scaled than leading commercial 5\,nm.'' Capability is demonstrated; TSMC-equivalent economics are not, and the metrics that would settle it are undisclosed [A]. But the strategic arithmetic is wafers times yield times time, and SMIC's advanced-node capacity climbs from about 45,000 to 80,000 wafers a month between 2025 and 2027~\cite{semianalysis-ascend} behind Huawei-linked ``shadow fabs'' and a bank of pre-ban dies. Below the leading edge the story is already solar-shaped: China supplies 37\% of global 22--40\,nm output and about 70\% of mature-node capacity expansion, which is where the EV and robotics stack lives.

Memory is the quiet revolution. CXMT scaled from about 70,000 wafers a month in 2022 to 200,000--300,000 by mid-2026, targeting 600,000, and will be within about 25,000 wafers a month of Micron by year-end---the world's fourth, and by trajectory third, DRAM producer~\cite{cxmt-capacity,citrini}, with DDR5 yields above 80\%. Its July 2026 Shanghai listing closed up 466\% at roughly \$460 billion~\cite{cxmt-ipo,cxmt-reuters}, which on under 7\% floated is a sentiment indicator and nothing more. The 2025--26 shortage handed CXMT the swing-supplier role and flipped the narrative from ``Chinese dumping'' to Chinese pricing power within twelve months---a moment in a panic market, not proof of durable pricing power [A]. HBM is the ledger's honest exception: CXMT is in production at HBM2/2E, three to four years behind SK hynix, disclosed no dedicated commercial HBM investment in its own prospectus, and is reported to start limited HBM3E production around the end of 2026 at 20--30\% yield against about 70\% at SK~hynix~\cite{hw-cxmt-hbm3e} [A]; no public teardown has yet found Chinese-fabbed DRAM in a shipping accelerator~\cite{hbm-teardown}, and Huawei's stockpile of about thirteen million foreign stacks is assessed as effectively exhausted. Domestic HBM is therefore the single binding physical constraint on the compute element and the most important open physical variable in the argument. The strategy is double-tracked---sprint on HBM three years behind, lead on the commodity LPDDR class the capacity-first architecture was designed to need---so China has not made HBM irrelevant but is engineering a position in which it ceases to be the \emph{sole} path to frontier-scale AI. Two 2026 disclosures cut both ways. Spot DRAM is up $7.1\times$ and NAND $6.9\times$ from the mid-2025 trough, with memory rising from 52\% to 63\% of AI-chip component spend~\cite{hc2026-memory-wall} [A]---which raises the cost of the West's HBM-centric architecture \emph{and} of the commodity-DRAM socket this thesis depends on, and the second effect is the one this book has to own. And ``HBM4'' has stopped specifying a part, quoted at 2.0 and at 3.0\,TB/s per stack by different vendors, with a customer-specific base die explicitly \emph{not} JEDEC-standardized~\cite{hc2026-hbm4-spread,hc2026-chbm} [P]---which relocates the frontier off the memory cell, where CXMT is measurably closing, and onto packaging and a custom logic die.

\emph{One month, two verdicts.} CXMT earned third-quarter net profit of RMB~54.4B on revenue of RMB~118.71B from commodity DRAM alone, while SK~hynix called its DRAM and HBM ``effectively sold out through 2029''~\cite{hw-cxmt-q3,hw-skhynix-q3} [V; P]: a self-financing cheap tier beside an expensive tier closed to outsiders strengthens the memory arbitrage, and the HBM3E yield strengthens the HBM concession. Against both, contract-DRAM growth slowed from $+$93--98\% in the first quarter to $+$13--18\% in the third~\cite{trendforce-dram} [A].

The tool base is the deepest dependency, and there the trend, not the level, is the message: Naura fifth in global equipment revenue, domestic tools 20--30\% of Chinese fab demand against 10\% three years earlier~\cite{toolmakers-top20,naura}; lithography the hard floor, with domestic EUV no earlier than about 2030; a US EDA ban rescinded within six weeks in exchange for rare-earth relief while domestic EDA share doubled; Big Fund III at RMB~344 billion~\cite{bigfund,semi-wfe}. Targets are missed and reset, but the direction never reverses, and each US control action functions as a demand guarantee for the domestic substitute. Ranking export controls by the effective delay they buy---EUV real years; HBM less than assumed, given stockpiles, diversion and the architectural detour; accelerator controls largely converted into a domestic market-reservation policy; EDA near zero---would have predicted, in 2022, most of what has since happened. The chokepoint strategy assumed China would rebuild the Western machine and could be stopped at its scarcest parts. China is building a different machine. Whether it works end to end---logic yield, packaging, memory qualification, software maturity, and capacity-first training performance demonstrated simultaneously in one frontier run---is the integration test not yet passed \full{13}.

\section{The watt: energy as the fourth element}
\label{c:energy}

Every token is electricity through silicon, and the AI war's longest-lead-time input is neither the model, the chip, nor the fab---it is the watt. Nameplate capacity added is not deliverable power at a computing site, so the ledger has four rows and China leads on all of them: about 540\,GW of new nameplate in 2025 against a record 53\,GW in the United States~\cite{nea-2025,us-additions}; 10,573\,TWh of generation against 4,430~\cite{ember-2026}; firm capacity from a record 95\,GW of coal commissioned as flexible backup, 66\,GW of nuclear heading toward 110 and two thirds of the world's grid storage, against American gas turbines booked to 2029--30, with GE Vernova's slots running to 2032; and deliverable headroom from 45 ultra-high-voltage lines with state-directed siting, against a US interconnection queue of roughly 2,000\,GW at a median wait above three years~\cite{lbl-queued}, transformers on 24--60-month lead times, and moratorium bills in at least twelve states. OpenAI's request for 100\,GW a year of new US capacity is nearly double the best year the American grid has ever had.

\begin{figure}[htbp]
\centering
\begin{tikzpicture}
\begin{axis}[bookbar, height=4.3cm, bar width=12pt,
  ylabel={Generating capacity added (GW)},
  title={Annual additions to generating capacity},
  symbolic x coords={2023,2024,2025},
  xtick=data, ymin=0, ymax=620,
  nodes near coords, nodes near coords style={font=\scriptsize\color{inkSec}},
  legend style={at={(0.02,0.98)}, anchor=north west},
]
\addplot[fill=sOne] coordinates {(2023,300) (2024,429) (2025,540)};
\addlegendentry{China}
\addplot[fill=sTwo] coordinates {(2023,32) (2024,45) (2025,53)};
\addlegendentry{United States}
\end{axis}
\end{tikzpicture}

\vspace{1mm}

\begin{tikzpicture}
\begin{axis}[bookbar, height=3.7cm, width=0.72\linewidth, bar width=15pt,
  ylabel={GW},
  title={The queue and the ask, for scale},
  symbolic x coords={{US 2025 built},{OpenAI's annual ask},{US interconnection queue}},
  xtick=data, ymin=0, ymax=2400,
  xticklabel style={font=\scriptsize\color{inkSec}, align=center},
  nodes near coords, nodes near coords style={font=\scriptsize\color{inkSec}},
]
\addplot[fill=sFour] coordinates
  {({US 2025 built},53) ({OpenAI's annual ask},100) ({US interconnection queue},2060)};
\end{axis}
\end{tikzpicture}
\caption[The energy asymmetry]{The fourth element. \emph{Above}: annual additions to generating capacity, China
against the United States---China added roughly ten times as much in 2025, into falling utilization hours, while
the US record year was 53\,GW [V]. \emph{Below}: the American constraint stated
three ways---what was built in the best year on record, the 100\,GW/year one AI laboratory told the White House
the country needs, and the $\sim$2,060\,GW sitting in the interconnection queue at median waits over three years
[V/P]. One side's binding constraint is capital; the other's is physics and
permitting, and permitting is the harder one to finance away.}
\label{fig:energy}
\end{figure}


The asymmetry is not merely quantitative; it is being aimed. ``East Data, West Computing'' placed eight national computing hubs where curtailed renewables are, and in November 2025 the western provinces cut power prices by up to 50\%, to about RMB~0.4 per kilowatt-hour, exclusively for facilities running on approved domestic chips, the same week Beijing barred foreign AI chips from state-funded data centers~\cite{china-dc-power,datacenter-mandate}. That single instrument welds three elements together: it converts grid surplus into a subsidy, the subsidy into demand for Ascend and Cambricon silicon, and the silicon into tokens priced below Western marginal cost, socializing part of the domestic chips' 30--50\% efficiency penalty. Magnitudes matter---electricity is only about 8\% of a US AI data center's three-year cost, and Chinese spot prices overlap American ones---so the advantage is not the tariff but the headroom: one side's constraint is money, the other's is physics and permitting, and permitting is the harder constraint to finance away. A busbar-level casebook across five American regions and five Chinese hubs finds the two systems failing at opposite ends of the same pipe: the United States on time to interconnect, China on the durability of administered tariffs. The latency objection to the western hubs is substantially wrong for now---network latency is about 5\% of a time-to-first-token budget, and Alibaba put an inference supernode into commercial service in Inner Mongolia in August 2026 serving Kimi K3 and Qwen3.8-Max worldwide---and Western capacity is underused because it was built for a pretraining boom and sits far from paying customers, which is an economic problem, not a physical one. The watt is where the precedent wars and the AI war become one war: solar, batteries and UHV transmission are the same commons Chapter~\ref{c:precedents} described, and the fourth element was the first three's prize \full{14}.


\emph{The siting question, with both sides at full strength.} The American ``moratorium'' is mostly a taxonomy problem. Exactly \textbf{one} statewide moratorium is in force---New York's executive order of 14~July 2026, reaching only \emph{discretionary state} environmental permits above 50\,MW and \textbf{not local permits at all}; Pennsylvania's order is a conditional regime, Loudoun County voted only to \emph{draft} one, Texas froze an interconnection queue, and \textbf{twelve of fourteen state ``moratoria'' are unpassed bills}~\cite{moratorium-taxonomy} [V/A]. On scale the two headline numbers measure different objects: \$130B of projects ``blocked or delayed'' (media scan, no megawatt figures, undisclosed funders) against \textbf{2.3\,GW actually delayed} of $\sim$20\,GW inside restricted areas, three projects accounting for nearly all of it, versus 38\,GW forecast for 2027 delivery~\cite{moratorium-scale} [P/A]. The strongest case \emph{for} restriction is that PJM billed \textbf{\$6.2B} of capacity cost against data centres \emph{forecast but not built}, while clearing below its own reliability requirement at the price cap---and Texas's queue holds 474\,GW against a $\sim$90\,GW peak~\cite{moratorium-attribution} [A]. The strongest case \emph{against} is that 76--126\,GW of headroom exists for loads accepting curtailment of a quarter to one per cent of annual energy, concentrated into 85--366 hours~\cite{flexible-load-headroom} [A]---thirty-three to fifty-five times what moratoria delay, needing no new generation---and that the best-identified causal study finds data centres \emph{lowered} retail rates 3.5\%, while a decomposition \textbf{commissioned by a consumer advocate} attributed \$8.7bn of PJM's \$12.5bn price spike to an accreditation rule change and two plant retirements rather than to load~\cite{moratorium-attribution} [A]. Much of the apparent contradiction dissolves on inspection: the studies measure capacity, wholesale and retail prices respectively, and fixed-cost spreading and marginal dispatch are both real and act on different bill components. Capacity is not fleeing abroad---no source establishes it; what moved is \emph{within} the US, Northern Virginia down 29\% while Atlanta became the most active market. And the federal government has \textbf{expressly declined to preempt} local siting, carving data centres out of its own AI-preemption order~\cite{federal-nopreempt,us-capacity-not-fleeing} [V].

\emph{And the Chinese comparison cuts both ways.} State siting authority bought genuine velocity---39 completed ultra-high-voltage lines, fast western interconnection---but \textbf{not cheap power}: in the \emph{coastal} provinces where Chinese data centres actually are, industrial rates sit near or above Virginia's. East Data West Computing delivered measured energy savings of \textbf{4--12\%} with negligible emissions or renewables benefit, and 75\% of digital-industry income remains in East China. Against a 60\% utilization target, \textbf{Beijing's own estimate is 20--30\%}, with \textbf{100+ state-backed projects cancelled in eighteen months}~\cite{china-dc-utilization} [A]. Velocity bought stranded assets, and the two are causally linked: building without a local veto or a demand test means building where the demand is not. Neither system has solved allocation \full{14}.

\emph{The autumn moved both ends of the pipe.} China's January--September additions fell about 31\% to 252.93\,GW, solar's share from 65.5\% to 43.8\%, while firm additions rose~\cite{cap-china-additions} [A; M on the percentage], and on 16~October the State Council directed that computing be aligned with power supply~\cite{cap-state-council-oct} [V]: an administered system now rations the match, on a schedule the American queue still cannot meet. FERC rejected PJM's curtailable large-load service on 9~October~\cite{cap-ferc-pjm-oct} [V], Oracle declared force majeure on its 2.5\,GW gas-fed Project Jupiter, and San Francisco, Oakland and San Diego adopted forty-five-day study moratoria, San Diego's extendable to nearly two years---little capacity, much signal~\cite{cap-dc-obstacles,cap-moratoria-autumn} [A] \full{14}.

\section{Data and feedback: the fifth element}
\label{c:data}

If open weights commoditize the model and value migrates to adjacent scarcities, the most durable is the one input openness cannot replicate by construction: data that only deployment generates. An open-weight release transfers a compressed function of its training distribution; it cannot transfer proprietary enterprise workflows, scientific-instrument streams, clinical data, skilled human correction at scale, or production failure traces---each generated by use, appropriated by whoever owns the deployment relationship, and scarce once the model is free. That explains why the American flywheel's rent recapture survives open weights---the enterprise-workflow row, which Chinese open models cannot reach from outside, is the row a \$47 billion enterprise run rate is built on---why China's embodied push is rational even if humanoids disappoint, and why evaluation and incident data is currently unowned, the same vacancy as the certification layer. Synthetic data is the sharpest challenge, and it amplifies whoever already holds the teacher or the verified anchor rather than substituting for either: distillation relocates the frontier's rent, it does not abolish the frontier. The fifth element is ground truth at scale with provenance---and that is what cannot be downloaded \full{15}.

\section{Embodied AI: the loop closes on the factory floor}
\label{c:embodied}

Everything above runs one way: China's manufacturing victories funded and shaped its digital-AI strategy. Embodied AI runs the loop the other way. ``Humanoid robot'' names a product, not an architecture: the United States leads the foundation and world-model layers, China leads motion, control, actuation and the deployment environment, and the contested middle---task policy---is decided by a physical-data flywheel whose scarce input is deployment surface, which cannot be scraped, only produced. China has built that loop's infrastructure as public works, in more than forty state-backed data-collection centers, and its component position is the precedent industries with new part numbers: roughly 70\% of Tesla's Optimus supply chain is in China and excluding it would triple the cost~\cite{optimus-china}. The American humanoid, like the American solar panel of 2010, is being designed atop the adversary's supply chain.

The scoreboard reads like a rout---more than 13,000 humanoids shipped in 2025, about 85\% by Chinese firms~\cite{fortune-humanoid}, Unitree profitable on 5,500 units~\cite{unitree-ipo}---and it deflates at once: shipments are not labor. Most Chinese units go to demonstrations and state data farms, the best-documented autonomous performances are American, and the metric that would settle the contest---useful work per robot-hour at fleet scale---is published by nobody. The decisive uncertainty is whether model capability or physical integration binds, and the author's forecast [S] is that integration wins in structured environments within the decade, which is China's terrain, while unstructured domains stay capability-bound well past 2030. If both close, the two flywheels fuse into one self-reinforcing industrial system---AI improving the factories that build the machines that generate the data that improves the AI---and the country that owns that loop owns something qualitatively beyond anything Chapter~\ref{c:precedents} described \full{16}.
